\section{Introduction}
\label{sec:introduction}

Intercity freight service combines dispersed collection and distribution with transportation between cities. Pooling shipments can reduce routing effort per shipment, but orders have release times and delivery deadlines. Vehicles and freight must also be available at the right place and time: a local collection wave creates a subsequent line-haul requirement, and a direct trip changes the fleet available in each city. The operational question is how to coordinate these services when local capacity itself depends on the shipments pooled into a wave.

We study finite-horizon scheduling for a two-city system with fixed fleets and given order time windows. Human-driven vehicles perform collection and distribution, while a separate fleet provides automated hub-to-hub transportation, as in autonomous transfer hub architectures \citep{Lee2024}. The flexible-service model also permits human-driven direct transportation and independent empty repositioning. The objective combines vehicle operating costs and penalties for unserved demand. A wave groups stops assigned to a common service interval, city, and shipment direction; its total quantity differs from the load carried by one vehicle.

Continuous routing approximations provide a way to retain this distinction without selecting individual customer tours. The square-root scaling of \citet{Beardwood1959}, spatial partitioning methods \citep{Karp1977}, and fleet-routing approximations \citep{Daganzo1984} motivate the local service-time model. Closely related, \citet{Banerjee2022} connect single-vehicle zone-area maximization to fleet sizing and service-region partitioning for same-day delivery. Their outputs are tactical design decisions. Here fleet sizes and order windows are inputs, and the outputs are dated services, freight flows, and vehicle movements. The shared methodological element is an aggregate routing approximation.

The paper makes two main contributions. First, it integrates wave-dependent local capacity into a finite-fleet scheduling model. The routing approximation yields capacity tables indexed by service duration and vehicle count, which enter a mixed-integer linear program through binary selectors. Cumulative vehicle inventories and freight precedence connect local waves to intercity services, direct trips, and empty movements. For a fixed local service arc, we show that modeled capacity per vehicle is nondecreasing and pooled capacity is superadditive. These properties explain the capacity mechanism used by the formulation.

Second, it provides a restricted analytical interpretation of wave and vehicle-load choices. In a stationary common-wave benchmark, we derive optimal collection and distribution loads conditional on the wave quantity, retaining both vehicle capacity and wave quantity as load bounds. Under constant arrivals, full-backlog dispatch, and separable batch-service costs, repeating a minimizer of the per-shipment objective is average-cost optimal among admissible successive wave sizes. The benchmark shares the routing primitive with the scheduling model but has continuous loads, a lead-time charge, and no resource competition between batches. It supplies a mechanism interpretation rather than an established relaxation or steady-state limit of the MILP. Complete QVI verification and a comparison on a common integer full-reset threshold domain are retained in the appendices.

Rolling execution supports the main scheduling model by carrying initiated tasks and inherited inventories across computational windows. Capacity construction, scheduling, and execution therefore form the main development of the paper; the stationary analysis is a separate explanatory step. Section~\ref{sec:related} reviews related work. Sections~\ref{sec:capacity} and~\ref{sec:model} present service capacities and finite-horizon scheduling with rolling implementation. Section~\ref{sec:analysis} develops the stationary benchmark, and Section~\ref{sec:discussion} concludes.
